% !PRINTED 421
\weberpage{421}

If we form from these invariant forms of the tetrahedron functions of $x_1:x_2$ and set $\vartheta = x_1:x_2$, we obtain functions which remain unchanged by the linear fractional tetrahedral substitutions, and which are therefore roots of resolvents of degree 5 of the icosahedral equation. The icosahedral equation itself is obtained in two forms, if one sets
\weq{10}{H^3 = zf^5,\quad T^2 = z_1f^5,\quad z + z_1 = 1728}
and the resolvents depend rationally on $z$.

If for the moment we denote by $\varphi(x_1, x_2)$ any absolutely invariant form of the tetrahedron, then this form passes, by any one of the substitutions $Q\Theta$, into
\weq{11}{\varphi_\varepsilon = \varphi(\varepsilon^{-2}x_1,\ \varepsilon^2x_2)}
and this function remains unchanged by the substitutions of the group $\Theta^{-1}Q\Theta$.

\gesperrt{From this it follows that every symmetric function of the five forms $\varphi_\varepsilon$ is an invariant of the icosahedral group, and can therefore, by the theorem of \S.\,61, be expressed by the fundamental forms $f, H, T$.}

This theorem permits a comparatively easy computation of the resolvents.

The simplest functions that we can use as roots of resolvents are
\weq{12}{r = \frac{t^2}{f},\quad u = \frac{f^2t}{T},\quad v = \frac{f\tilde\omega}{H}.}

If we first observe that the symmetric fundamental functions of the five functions $t_\varepsilon$ are icosahedral invariants of the degrees $6, 12, 18, 30$, then, since the degree numbers $6, 18$ do not occur among these invariants, there follows from the theorems in \S.\,61 an equation of the form
\weq{13}{t^5 + aft^3 + bf^2t + cT = 0,}
in which $a, b, c$ are numerical coefficients. To determine them, we denote by $S(\varphi)$ the sum of the five functions $\varphi_\varepsilon$ in (11) and by $\Pi(\varphi)$ their product, and obtain from the \gesperrt{Newton} formulae (Bd.\,I, \S.\,42):
\weq{14}{af = -\tfrac{1}{2}S(t^2),\quad bf^2 = \tfrac{1}{2}a^2f^2 - \tfrac{1}{4}S(t^4),\quad cT = -\Pi(t).}

In the sums $S(\varphi)$ only those terms $x_1^h x_2^k$ have a coefficient different from zero in which $h - k$ is divisible by 5. One therefore obtains $a, b, c$ by equating the
%CONT>>
